\documentclass[letterpaper,10pt]{article}

\usepackage[margin=0.75in]{geometry}
\usepackage{times}
\usepackage{helvet}
\usepackage{courier}
\usepackage{microtype}
\usepackage{booktabs}
\usepackage{array}
\usepackage{graphicx}
\usepackage{url}
\usepackage{xcolor}
\usepackage{caption}

\setlength{\pdfpagewidth}{8.5in}
\setlength{\pdfpageheight}{11in}
\emergencystretch=4em
\urlstyle{rm}
\def\UrlFont{\rm}
\makeatletter
\g@addto@macro{\UrlBreaks}{\do\/\do\_\do\-}
\makeatother
\frenchspacing

\newcommand{\safe}{\textsc{Safe}}
\newcommand{\cont}{\textsc{Continue}}
\newcommand{\abstain}{\textsc{Abstain}}

\title{Supplementary Material for\\
Local Evidence Is Not Completion:\\
Evidence-Condition Geometry for Bounded Completion-Audit Workflows}
\author{Anonymous Submission}
\date{}

\begin{document}
\sloppy
\maketitle

\section*{Scope of the Supplement}

This anonymous supplementary material contains reproducibility and audit details
that are intentionally kept out of the main technical-content pages. These
tables support the evidence-condition controller evaluation but are not used as
additional main effectiveness evidence. In particular, the workflow pilot is a
workflow-shape/interface validation only, and the scalar singleton/Chao proxy is
a negative control because it does not encode source-route evidence conditions.

\section*{Reproducibility Notes}

\paragraph{Script entry points.}
All paths are relative to the repository root. The supplement and paper figures
are assembled from \path{2/analysis/research_object_geometry/real_agent_pilot/}
by running the Python entry point
\path{credibility_supplement/run_credibility_supplement.py}. The frozen task logs
used by this assembler are produced by
\path{scripts/run_blind_policy_task.py},
\path{scripts/run_blind_code_task.py},
\path{external_validation_requests/run_external_requests_validation.py},
\path{controller_validation_v1/run_controller_validation_v1.py},
\path{controller_validation_v1/run_controller_validation_v2.py}, and
\path{external_validation_v2/run_external_validation_v2.py}. The controlled
stress tests are in \path{2/analysis/research_object_geometry/}:
\path{controlled_geometry_simulation_v2.py} and
\path{controlled_exposure_localization_v3.py}.

\paragraph{Seeds and thresholds.}
Generated-task method validation, external repository repair validation, and
external boundary validation use 200 fixed seeds per challenger or controller
policy before oracle scoring. Safe-state validation uses 6000 order/budget
perturbations. The controller-validation thresholds are
\(\tau_s=0.75\) for source-route support and \(\tau_g=0.70\) for exposure
Gini. The 0.90 bounded-oracle recall threshold is used only for post-hoc
false-certification evaluation and is not visible to the controller.

\paragraph{Data and log format.}
Primary run logs are JSONL action/evidence ledgers. Each event records a task
identifier, source family, source-route stratum, route, action or scan record,
discovered item identifier when present, new-item flag, and cost fields used for
post-stop repair or confirmation-audit accounting. Aggregated CSVs in
\texttt{credibility\_supplement/results/} contain condition summaries,
controller decisions, repair confidence intervals, seeded safe-state summaries,
threshold/budget sweeps, oracle summaries, route-pattern examples, and leakage
controls.

\paragraph{Oracle CSV and JSONL fields.}
Oracle item files use the fields
\texttt{task\_id}, \texttt{item\_id}, \texttt{oracle\_label},
\texttt{oracle\_bucket}, \texttt{source\_path}, \texttt{source\_family},
\texttt{source\_route\_stratum}, \texttt{reportable}, and, for external
repository JSONL files, a representative \texttt{line}. Source-route count CSVs
use \texttt{source\_family}, \texttt{oracle\_bucket}, and
\texttt{oracle\_items}. Oracle labels and oracle totals are loaded only after
trajectories, repair orderings, and challenger choices are fixed.

\paragraph{Figure generation.}
The same supplement assembler writes
\path{main_results_overview.png}, \path{controller_decision_matrix.png}, and
\path{repair_sensitivity_summary.png} to \path{2/paper/figures/}. Run the
assembler above, then compile from \path{2/paper/} with \texttt{latexmk -pdf
-interaction=nonstopmode -outdir=build} and either
\path{evidence_condition_geometry_aaai_v2.tex} or
\path{supplementary_material.tex}.

\paragraph{Computing environment.}
The checked environment is Windows 11 with Python 3.13.9, NumPy 2.3.5, pandas
2.3.3, and Matplotlib 3.10.6. The code uses only local frozen snapshots for the
external repository audits; no network access is required to reproduce the
reported tables and figures from the included logs and snapshots.

\begin{table}[t]
\centering
\small
\setlength{\tabcolsep}{3.2pt}
\caption{Boundary case showing why eligibility is not proof. The state passes
the source-route eligibility gate, but post-hoc recall remains below the
completion threshold; repair-aware control therefore returns CONTINUE rather
than SAFE.}
\label{tab:eligibility_boundary}
\resizebox{\columnwidth}{!}{%
\begin{tabular}{llrrrrll}
\toprule
Task & State & Support & Gini & Recall & Residual+ & Elig.-only & Full \\
\midrule
urllib3 & route\_partitioned & 0.800 & 0.647 & 0.835 & yes & SAFE & CONTINUE \\
\bottomrule
\end{tabular}}
\end{table}

\begin{table}[t]
\centering
\small
\setlength{\tabcolsep}{2.8pt}
\caption{Seeded unsafe/complete controller decision counts for
Figure~\ref{fig:controller_variants}. Unsafe and Complete are the seeded
denominators; FCR is SAFE-on-unsafe divided by Unsafe, and safe coverage is
SAFE-on-complete divided by Complete. These denominators differ from the
one-state eligibility boundary in Table~\ref{tab:eligibility_boundary}. Oracle
labels are used only after decisions are fixed. S/C/A reports
SAFE/CONTINUE/ABSTAIN counts.}
\label{tab:main_controller_counts}
\resizebox{\columnwidth}{!}{%
\begin{tabular}{lrrrrrr}
\toprule
Decision rule & Unsafe & Complete & Unsafe S/C/A & Complete S/C/A & FCR & Safe cov. \\
\midrule
Naive stop & 400 & 1200 & 400/0/0 & 1200/0/0 & 1.000 & 1.000 \\
Source-only & 400 & 1200 & 400/0/0 & 1200/0/0 & 1.000 & 1.000 \\
Verifier-gate & 400 & 1200 & 0/0/400 & 1200/0/0 & 0.000 & 1.000 \\
Eligibility-only & 400 & 1200 & 0/0/400 & 1200/0/0 & 0.000 & 1.000 \\
Full & 400 & 1200 & 0/400/0 & 1200/0/0 & 0.000 & 1.000 \\
\bottomrule
\end{tabular}}
\end{table}

\begin{table*}[t]
\centering
\caption{Controller decision accounting. Oracle-safe and oracle-unsafe counts
are post-hoc denominators used only for evaluation. Cost is measured repair
cost, counted after a fixed repair target set is selected.}
\label{tab:controller_decisions}
\resizebox{\textwidth}{!}{%
\begin{tabular}{llrrrrrrrrrrrr}
\toprule
Task group & Evaluation set & \(n\) & Oracle-safe & Oracle-unsafe
& \#SAFE & \#CONTINUE & \#ABSTAIN & False SAFE & FCR
& Safe cov. & Abstain & Repair gain & Repair cost \\
\midrule
all tasks & observed stop states & 9 & 4 & 5 & 4 & 5 & 0 & 0 & 0.000 & 1.000 & 0.000 & 0.000 & 0.0 \\
external repos & observed external states & 5 & 2 & 3 & 2 & 3 & 0 & 0 & 0.000 & 1.000 & 0.000 & 0.000 & 0.0 \\
requests & seeded repairs & 1200 & 0 & 1200 & 0 & 1199 & 1 & 0 & 0.000 & -- & 0.001 & 122.992 & 2854.6 \\
urllib3 & seeded repairs & 1000 & 0 & 1000 & 0 & 998 & 2 & 0 & 0.000 & -- & 0.002 & 204.987 & 4831.8 \\
external repos & seeded repairs & 2200 & 0 & 2200 & 0 & 2197 & 3 & 0 & 0.000 & -- & 0.001 & 160.262 & 3753.3 \\
requests & seeded safe states & 3000 & 3000 & 0 & 3000 & 0 & 0 & 0 & 0.000 & 1.000 & 0.000 & 0.000 & 3252.5 \\
urllib3 & seeded safe states & 3000 & 3000 & 0 & 3000 & 0 & 0 & 0 & 0.000 & 1.000 & 0.000 & 0.000 & 3933.7 \\
external repos & seeded safe states & 6000 & 6000 & 0 & 6000 & 0 & 0 & 0 & 0.000 & 1.000 & 0.000 & 0.000 & 3593.1 \\
\bottomrule
\end{tabular}}
\end{table*}

\begin{table}[t]
\centering
\small
\setlength{\tabcolsep}{2.5pt}
\caption{Small claim-verification pilot with human-style audit routes. Routes
are support, contradiction, exception-path, config-default, and scope-boundary.
This deterministic pilot is a scoped sanity check for completion control, not a
human-annotated benchmark. S/C/A reports SAFE/CONTINUE/ABSTAIN counts.}
\label{tab:claim_verification_pilot}
\resizebox{\columnwidth}{!}{%
\begin{tabular}{lrrrrrrr}
\toprule
Policy & Unsafe & Safe & Unsafe S/C/A & Safe S/C/A & FCR & Safe cov. & Cost \\
\midrule
Naive stop & 4 & 2 & 4/0/0 & 2/0/0 & 1.000 & 1.000 & 0.0 \\
Source-only & 4 & 2 & 4/0/0 & 2/0/0 & 1.000 & 1.000 & 0.0 \\
Verifier-gate & 4 & 2 & 0/0/4 & 1/0/1 & 0.000 & 0.500 & 0.0 \\
Eligibility-only & 4 & 2 & 1/0/3 & 2/0/0 & 0.250 & 1.000 & 0.0 \\
Full controller & 4 & 2 & 0/4/0 & 2/0/0 & 0.000 & 1.000 & 3.5 \\
\bottomrule
\end{tabular}}
\end{table}

\begin{table*}[t]
\centering
\caption{Per-task decision breakdown on observed stop states. FCR is
SAFE-on-unsafe divided by oracle-unsafe stop states for that task; safe coverage
is SAFE-on-safe divided by oracle-safe stop states for that task. Oracle labels
are post-hoc only.}
\label{tab:per_task_decision_breakdown}
\resizebox{\textwidth}{!}{%
\begin{tabular}{llrrrrrrr}
\toprule
Task & Policy & Unsafe & Safe & SAFE-on-unsafe & FCR & Safe cov. & CONTINUE & ABSTAIN \\
\midrule
policy-docset & Naive stop & 1 & 1 & 1 & 1.000 & 1.000 & 0.000 & 0.000 \\
policy-docset & Source-only & 1 & 1 & 1 & 1.000 & 1.000 & 0.000 & 0.000 \\
policy-docset & Verifier-gate & 1 & 1 & 1 & 1.000 & 1.000 & 0.000 & 0.000 \\
policy-docset & Eligibility-only & 1 & 1 & 0 & 0.000 & 1.000 & 0.000 & 0.500 \\
policy-docset & Full controller & 1 & 1 & 0 & 0.000 & 1.000 & 0.000 & 0.500 \\
code-repo & Naive stop & 1 & 1 & 1 & 1.000 & 1.000 & 0.000 & 0.000 \\
code-repo & Source-only & 1 & 1 & 1 & 1.000 & 1.000 & 0.000 & 0.000 \\
code-repo & Verifier-gate & 1 & 1 & 1 & 1.000 & 1.000 & 0.000 & 0.000 \\
code-repo & Eligibility-only & 1 & 1 & 0 & 0.000 & 1.000 & 0.000 & 0.500 \\
code-repo & Full controller & 1 & 1 & 0 & 0.000 & 1.000 & 0.000 & 0.500 \\
requests & Naive stop & 1 & 1 & 1 & 1.000 & 1.000 & 0.000 & 0.000 \\
requests & Source-only & 1 & 1 & 1 & 1.000 & 1.000 & 0.000 & 0.000 \\
requests & Verifier-gate & 1 & 1 & 1 & 1.000 & 1.000 & 0.000 & 0.000 \\
requests & Eligibility-only & 1 & 1 & 0 & 0.000 & 1.000 & 0.000 & 0.500 \\
requests & Full controller & 1 & 1 & 0 & 0.000 & 1.000 & 0.000 & 0.500 \\
urllib3 & Naive stop & 2 & 1 & 2 & 1.000 & 1.000 & 0.000 & 0.000 \\
urllib3 & Source-only & 2 & 1 & 2 & 1.000 & 1.000 & 0.000 & 0.000 \\
urllib3 & Verifier-gate & 2 & 1 & 2 & 1.000 & 1.000 & 0.000 & 0.000 \\
urllib3 & Eligibility-only & 2 & 1 & 1 & 0.500 & 1.000 & 0.000 & 0.333 \\
urllib3 & Full controller & 2 & 1 & 0 & 0.000 & 1.000 & 0.667 & 0.000 \\
\bottomrule
\end{tabular}}
\end{table*}

\begin{table*}[t]
\centering
\caption{Lightweight external baselines and sanity checks. These are not new
task suites: source-only is a direct stopping baseline, random/high-potential
are repair-target baselines under the same budget, and the scalar singleton/Chao
proxy is a negative control that ignores source-route geometry.}
\label{tab:lightweight_baselines}
\resizebox{\textwidth}{!}{%
\begin{tabular}{lllll}
\toprule
Task & Baseline & Measurement & Value & Note \\
\midrule
requests & Source-only stop & FCR if accepted & 1.000 & accepts localized route evidence \\
requests & random expansion & mean repair gain & 45.5 & same budget, post-stop repair \\
requests & high-potential expansion & mean repair gain & 177.0 & same budget, post-stop repair \\
requests & residual-potential expansion & mean repair gain & 177.0 & same budget, post-stop repair \\
urllib3 & Source-only stop & FCR if accepted & 1.000 & accepts localized route evidence \\
urllib3 & random expansion & mean repair gain & 92.5 & same budget, post-stop repair \\
urllib3 & high-potential expansion & mean repair gain & 275.0 & same budget, post-stop repair \\
urllib3 & residual-potential expansion & mean repair gain & 329.0 & same budget, post-stop repair \\
requests & Singleton/Chao scalar stop & would stop? & no & Chao1=496.0 \\
urllib3 & Singleton/Chao scalar stop & would stop? & no & Chao1=9180.0 \\
\bottomrule
\end{tabular}}
\end{table*}

\begin{table*}[t]
\centering
\caption{Negative control: scalar singleton/Chao proxy on homogeneous runs. The
proxy observes item-count sparsity but does not encode source-route condition,
so it is not a competitive controller baseline for certificate mismatch.}
\label{tab:chao_proxy}
\resizebox{\textwidth}{!}{%
\begin{tabular}{lrrrrrrrll}
\toprule
Task & Observed & Oracle total & Recall & \(f_1\) & \(f_2\)
& \(f_1/\mathrm{obs}\) & Chao1 & Scalar stop & False if stop \\
\midrule
\texttt{policy\_docset\_v1} & 17 & 24 & 0.708 & 17 & 0 & 1.000 & 153.0 & No & No \\
\texttt{code\_repo\_v1} & 6 & 20 & 0.300 & 6 & 0 & 1.000 & 21.0 & No & No \\
\texttt{requests} & 31 & 298 & 0.104 & 31 & 0 & 1.000 & 496.0 & No & No \\
\texttt{urllib3} & 135 & 699 & 0.193 & 135 & 0 & 1.000 & 9180.0 & No & No \\
\bottomrule
\end{tabular}}
\end{table*}

\begin{table}[t]
\centering
\caption{Sensitivity sweep numeric summary. Ranges aggregate threshold settings
or repair-budget settings within each external task.}
\label{tab:sensitivity_summary}
\resizebox{\columnwidth}{!}{%
\begin{tabular}{llllllr}
\toprule
Task & Sweep & SAFE rate & CONTINUE rate & ABSTAIN rate & Max FCR & Mean repair cost \\
\midrule
\texttt{requests} & threshold & 0.000-0.000 & 0.995-1.000 & 0.000-0.005 & 0.000 & 2854.6 \\
\texttt{urllib3} & threshold & 0.000-0.000 & 0.990-1.000 & 0.000-0.010 & 0.000 & 4831.8 \\
\texttt{requests} & budget & 0.000-0.000 & 0.635-1.000 & 0.000-0.365 & 0.000 & 2966.4 \\
\texttt{urllib3} & budget & 0.000-0.000 & 0.640-1.000 & 0.000-0.360 & 0.000 & 4036.5 \\
\bottomrule
\end{tabular}}
\end{table}

\begin{table}[t]
\centering
\caption{External oracle construction summary. Pattern-defined external oracles
are line-level source-route evidence occurrences from frozen local snapshots.}
\label{tab:oracle_appendix_summary}
\resizebox{\columnwidth}{!}{%
\begin{tabular}{lrrrrl}
\toprule
Repo & Version & Items & Routes & Nonempty strata & Granularity \\
\midrule
\texttt{requests} & 2.32.5 & 298 & 4 & 18 & line-level source-route \\
\texttt{urllib3} & 2.5.0 & 699 & 5 & 24 & line-level source-route \\
\bottomrule
\end{tabular}}
\end{table}

\begin{table*}[t]
\centering
\caption{Small logged workflow pilot for workflow-shape/interface validation
only. The pilot fixes the prompt and records independent contexts,
action/evidence events, localized stop proposals, and the controller-facing
decision record; it is not used as main effectiveness evidence.}
\label{tab:workflow_pilot}
\resizebox{\textwidth}{!}{%
\begin{tabular}{llrrrrrl}
\toprule
Task & Model & Agents & Independent contexts & Action events & Evidence events
& Stop proposals & Controller decision \\
\midrule
\texttt{T\_doc\_dynamic\_workflow\_smoke} & logged LLM-agent workflow & 3 & 3 & 39 & 35 & 3 & not scored; workflow-shape validation \\
\bottomrule
\end{tabular}}
\end{table*}

\begin{table*}[t]
\centering
\caption{Seeded safe-state validation. These are oracle-safe broad-complete states
with seeded order perturbations; the controller should remain SAFE rather than
over-reacting to extra repair opportunity.}
\label{tab:safe_state_validation}
\resizebox{\textwidth}{!}{%
\begin{tabular}{lllrrrrrrrrl}
\toprule
Task & State & Order/budget sweep & Runs & \#SAFE & \#CONTINUE & \#ABSTAIN & FCR & Safe cov. & Repair gain & Mean cost & Cost range \\
\midrule
\texttt{requests} & route\_partitioned & 5 orders; budgets 1-8 & 3000 & 3000 & 0 & 0 & 0.000 & 1.000 & 0.0 & 3252.5 & 157-7647 \\
\texttt{urllib3} & extended\_audit & 5 orders; budgets 1-8 & 3000 & 3000 & 0 & 0 & 0.000 & 1.000 & 0.0 & 3933.7 & 275-9072 \\
\bottomrule
\end{tabular}}
\end{table*}

\begin{table}[t]
\centering
\caption{Route-match consistency sanity check for external oracle positives. A
small sample of line-level positives was checked for route consistency;
ambiguous lines are retained as such rather than treated as full validation.}
\label{tab:oracle_sanity}
\resizebox{\columnwidth}{!}{%
\begin{tabular}{llrrrrl}
\toprule
Repo & Route & Sample & Positive & Precision & Ambiguous & Notes \\
\midrule
\texttt{requests} & compat\_route & 8 & 8 & 1.000 & 0 & line-level oracle-positive sample \\
\texttt{requests} & exception\_route & 8 & 8 & 1.000 & 0 & line-level oracle-positive sample \\
\texttt{requests} & timeout\_route & 8 & 8 & 1.000 & 0 & line-level oracle-positive sample \\
\texttt{requests} & tls\_route & 8 & 8 & 1.000 & 0 & line-level oracle-positive sample \\
\texttt{urllib3} & cleanup\_route & 8 & 7 & 0.875 & 1 & line-level oracle-positive sample \\
\texttt{urllib3} & exception\_route & 8 & 8 & 1.000 & 0 & line-level oracle-positive sample \\
\texttt{urllib3} & retry\_route & 8 & 8 & 1.000 & 0 & line-level oracle-positive sample \\
\texttt{urllib3} & timeout\_route & 8 & 8 & 1.000 & 0 & line-level oracle-positive sample \\
\texttt{urllib3} & tls\_route & 8 & 8 & 1.000 & 0 & line-level oracle-positive sample \\
\bottomrule
\end{tabular}}
\end{table}

\begin{table}[t]
\centering
\caption{Leakage control for controller and repair decisions. Runtime decisions
use only visible evidence. Oracle labels, oracle totals, post-hoc recall, and
undiscovered true item counts are used only after trajectories and challenger
choices are fixed.}
\label{tab:leakage}
\resizebox{\columnwidth}{!}{%
\begin{tabular}{lll}
\toprule
Component & Runtime allowed & Oracle forbidden \\
\midrule
Exposure & visits, scans, actions & missing mass \\
Support/Gini & exposure counts & post-hoc recall \\
Potential & text, routes, matches & hidden true items \\
Decision & condition and repair & target distribution \\
\bottomrule
\end{tabular}}
\end{table}

\begin{table*}[t]
\centering
\caption{External oracle appendix detail. Each route uses a frozen regex pattern;
items are deduplicated by source, route, and line. To avoid unreadable regex
truncation, the appendix reports a short pattern category and one representative
line-level item per route; full regex rules and items are in the generated CSVs.}
\label{tab:oracle_route_examples}
\scriptsize
\resizebox{\textwidth}{!}{%
\begin{tabular}{llrllp{0.46\textwidth}}
\toprule
Repo & Route & Items & Pattern category & Dedup & Representative item \\
\midrule
\texttt{requests} & \texttt{tls\_route} & 102 & TLS/certificate/verification terms & source+route+line & adapters:tls\_route:103: \# cert path \\
\texttt{requests} & \texttt{timeout\_route} & 31 & timeout/connect/read-timeout terms & source+route+line & adapters:timeout\_route:120: self, request, stream=False, timeout=None, verify=True, cert=None, proxies=None \\
\texttt{requests} & \texttt{exception\_route} & 121 & exception/raise/error terms & source+route+line & adapters:exception\_route:136: raise NotImplementedError \\
\texttt{requests} & \texttt{compat\_route} & 44 & compatibility/deprecation terms & source+route+line & adapters:compat\_route:32: from .compat import basestring, urlparse \\
\texttt{urllib3} & \texttt{timeout\_route} & 135 & timeout/connect/read-timeout terms & source+route+line & connection:timeout\_route:137: timeout: \_TYPE\_TIMEOUT = \_DEFAULT\_TIMEOUT, \\
\texttt{urllib3} & \texttt{retry\_route} & 168 & retry/backoff/status terms & source+route+line & connectionpool:retry\_route:1011: retries, \\
\texttt{urllib3} & \texttt{tls\_route} & 111 & TLS/certificate/verification terms & source+route+line & connection:tls\_route:1005: cert, \\
\texttt{urllib3} & \texttt{exception\_route} & 170 & exception/raise/error terms & source+route+line & connection:exception\_route:1015: except BaseException: \\
\texttt{urllib3} & \texttt{cleanup\_route} & 115 & close/release/finally terms & source+route+line & connection:cleanup\_route:1016: ssl\_sock.close() \\
\bottomrule
\end{tabular}}
\end{table*}



\end{document}
